\subsection{Objective}

Task 5 requires determining the outline(s) of the characters in the image and
displaying them.  The input is the binary image selected in Task~4
(\texttt{task4\_binary.png}), in which the ten characters appear as bright
strokes on a dark background.  Because the binary image already separates the
characters from the background, outline detection reduces to following the
border between foreground and background for every character-sized connected
component.

\subsection{Contour Following Method}

The outlines were traced with Moore-neighbor contour following~\cite{contourtracing}, implemented
from first principles in \texttt{common.py}.  For one 8-connected foreground
component the algorithm works as follows.

\begin{enumerate}
  \item Start at the top-leftmost foreground pixel $b_0$; the backtrace pixel
  $c_0$ is its western neighbour, which is background by construction.
  \item Examine the eight neighbours of the current boundary pixel $b_k$ in
  clockwise order, beginning just after the backtrace pixel $c_k$, until the
  first foreground neighbour is found; it becomes $b_{k+1}$, and the
  neighbour examined just before it becomes the new backtrace pixel
  $c_{k+1}$.
  \item Repeat until the start pixel is re-entered from the same direction
  (Jacob's stopping criterion), which closes the contour.
\end{enumerate}

The walk is clockwise around the component, so the outer boundary is visited
in order and its pixels can be encoded as an 8-direction Freeman chain code,
with direction $0$ east, $1$ north-east, $2$ north, $3$ north-west, $4$ west,
$5$ south-west, $6$ south, and $7$ south-east.  Characters such as
\texttt{D}, \texttt{4}, \texttt{8}, \texttt{0}, and \texttt{A} also contain
enclosed holes.  The holes are the background regions inside a component that
do not touch the image border; each hole is identified by labelling the
inverted component and traced with the same follower, so the method returns
the outer outline \emph{and} every inner outline.

Tracing is performed on a padded crop of each character.  Before plotting
or exporting, a crop coordinate $(r,c)$ is converted back to the grayscale
sub-image coordinate $(r-p+y_0,c-p+x_0)$, where $p=1$ is the padding and
$(y_0,x_0)$ is the crop origin.  This places every outer and inner contour
on its own character.  The script verifies that all trace coordinates lie
inside the corresponding bounding box, that outer trace pixels are
foreground, and that hole trace pixels are background.

\subsection{Results}

\cref{fig:task5-overlay} overlays the traced outlines on the grayscale
sub-image: solid curves are outer boundaries and dashed curves are hole
boundaries, with one colour per character.  \cref{fig:task5-mask} shows the
boundary pixels alone, and \cref{fig:task5-example} zooms into the second
\texttt{4}, marking the trace start and the clockwise walking direction.

\begin{figure}[!htbp]
  \centering
  \includegraphics[width=\textwidth]{task5_outlines.png}
  \caption{Traced outlines overlaid on the grayscale sub-image.  Solid lines
  are outer boundaries; dashed lines are hole boundaries.}
  \label{fig:task5-overlay}
\end{figure}

\begin{figure}[!htbp]
  \centering
  \begin{subfigure}{0.66\textwidth}
    \centering
    \includegraphics[width=\textwidth]{task5_outline_mask.png}
    \caption{Boundary pixels only.}
    \label{fig:task5-mask}
  \end{subfigure}
  \hfill
  \begin{subfigure}{0.32\textwidth}
    \centering
    \includegraphics[width=\textwidth]{task5_example_trace.png}
    \caption{Zoom on the second \texttt{4}: start point and trace direction.}
    \label{fig:task5-example}
  \end{subfigure}
  \caption{Outline detection outputs.}
  \label{fig:task5-figures}
\end{figure}

\begin{table}[!htbp]
  \centering
  \caption{Per-character outline statistics.  ``Expected'' is the ground-truth
  glyph in \texttt{HD44780A00}.}
  \label{tab:task5-outlines}
  \begin{tabular}{clrrl}
    \toprule
    Order & Expected & Outer length & Holes & Hole lengths \\
    \midrule
    1 & \texttt{H} & 619 & 0 & -- \\
    2 & \texttt{D} & 379 & 1 & 270 \\
    3 & \texttt{4} & 368 & 1 & 164 \\
    4 & \texttt{4} & 340 & 1 & 172 \\
    5 & \texttt{7} & 408 & 0 & -- \\
    6 & \texttt{8} & 515 & 1 & 135 \\
    7 & \texttt{0} & 329 & 3 & 1, 2, 236 \\
    8 & \texttt{A} & 364 & 1 & 130 \\
    9 & \texttt{0} & 332 & 2 & 229, 3 \\
    10 & \texttt{0} & 325 & 1 & 230 \\
    \bottomrule
  \end{tabular}
\end{table}

\subsection{Analysis}

The follower returns a closed outer boundary for all ten characters; the total
outer boundary length of the row is $3{,}979$ pixels.  The hole counts agree
with the topology of most glyphs: \texttt{H} and \texttt{7} have no holes, and
\texttt{D}, \texttt{4}, and \texttt{A} each have one enclosed region.

Three details show where the binarization stage leaves its fingerprint on the
outlines.  First, the \texttt{8} yields one hole instead of the two expected
for its two loops: in the thresholded image a stroke gap opens the bottom
loop interior to the background, so only the top loop remains enclosed.
Second, the first \texttt{0} reports two tiny holes of $1$--$2$ pixels beside
its main hole, and the middle \texttt{0} one of $3$ pixels: these are spurious
holes created by small stroke gaps, not real openings of the glyphs.  Third,
all ten retained components still admit closed outer contours.  This does
not mean that the thresholding defects leave the glyph shapes unchanged:
the traced boundaries describe the selected binary mask, and the missing
loop of \texttt{8} changes its topology.  Hole counts therefore require
inspection before they are used as recognition features.

The JSON file \path{task_5/outputs/task5_outlines.json} stores the complete
ordered outer and hole traces, their starting pixels, and their full Freeman
chain codes.  All trace and start coordinates use $(\mathrm{row},
\mathrm{column})$ in the Task~3 sub-image.  A non-degenerate chain includes
the final step back to its start; a single-pixel hole is represented by its
start and an empty chain.  These data reproduce the boundary-pixel mask
without re-running the follower.  The earlier chain-code prefix and
statistics remain available as compact summaries.

\FloatBarrier
